% BertSum抽取结构：每句前加[CLS]、后加[SEP]，分段嵌入按句子奇偶交替；
% 各[CLS]顶层向量经句间Transformer层与逻辑回归给出入选概率。
\begin{tikzpicture}[
  x=1mm,y=1mm,figure picture,
  tok/.style={figure node,minimum width=10mm,minimum height=6mm,
    inner sep=1pt,draw=FigureInput,fill=FigureInputFill,font=\figurefont\scriptsize},
  seg/.style={font=\figurefont\scriptsize,text=FigureMuted},
  enc/.style={figure node,minimum height=7mm,inner sep=1.5pt,
    draw=FigureModel,fill=FigureModelFill,font=\figurefont\scriptsize},
  vec/.style={font=\figurefont\scriptsize},
  outc/.style={figure node,minimum width=12mm,minimum height=6mm,
    inner sep=1pt,draw=FigureOutput,fill=FigureOutputFill,font=\figurefont\scriptsize},
  row label/.style={font=\figurefont\scriptsize,text=FigureMuted,anchor=east},
  flow/.style={figure flow}
]
  % 三个句子，每句三个位置：[CLS]、句子词元、[SEP]
  \foreach \k/\s/\e in {1/{$s_1$词元}/E_A,2/{$s_2$词元}/E_B,3/{$s_3$词元}/E_A}{
    \pgfmathsetmacro{\xo}{(\k-1)*36}
    \node[tok] (c\k) at (\xo,0) {[CLS]};
    \node[tok,minimum width=12mm] (w\k) at (\xo+12.5,0) {\s};
    \node[tok] (p\k) at (\xo+25,0) {[SEP]};
    \node[seg] at (\xo+12.5,-6) {分段$\e$};
    \draw[flow] (c\k) -- (c\k |- 0,8.5);
    \draw[flow] (w\k) -- (w\k |- 0,8.5);
    \draw[flow] (p\k) -- (p\k |- 0,8.5);
    \node[vec] (t\k) at (\xo,23) {$\bm t_{\k}$};
    \draw[flow] (\xo,15.5) -- (t\k);
    \node[outc] (z\k) at (\xo,48) {$\hat z_{\k}$};
  }
  \node[enc,minimum width=108mm] (bert) at (48.5,12)
    {BERT编码器（词元嵌入＋分段嵌入＋位置嵌入）};
  \node[enc,minimum width=84mm] (inter) at (36,35)
    {句间Transformer层（$L=2$），输入加句子位置嵌入};
  \foreach \k in {1,2,3}{
    \pgfmathsetmacro{\xo}{(\k-1)*36}
    \draw[flow] (t\k) -- (\xo,31.5);
    \draw[flow] (\xo,38.5) -- (z\k);
  }
  \node[row label] at (-7,0) {输入};
  \node[row label] at (-7,23) {句向量};
  \node[row label] at (-7,48) {$\sigma(\bm w_o^{\top}\bm h_i+b_o)$};
\end{tikzpicture}
